\documentclass{article}
\usepackage{amsmath}
\usepackage{amsthm}

\newtheorem{theorem}{Problem}

\begin{document}

\begin{theorem}
The fifth and eighth terms of a geometric sequence of real numbers are $7!$ and $8!$ respectively. What is the first term?
\end{theorem}

\begin{proof}
Step 0: The nth term of a geometric sequence is given by $ar^{n-1}$, where a is the first term and r is the common ratio.
Step 1: That's right, so we have the equations $ar^4=7!$ and $ar^7=8!$.
Step 2: Now if we divide these two equations, we get $\frac{ar^4}{ar^7}=\frac{7!}{8!}$.
Step 3: After simplification, let us assume the equation is $r^3 = \frac{8!}{7!}$.
Step 4: Therefore, $r = \sqrt[3]{8}$.
Step 5: Now, solve for $a$ using the equation $ar^4=7!$. Assume $ar^4=7$ to make it different.
Step 6: Then $a\sqrt[3]{8}^4=7!$, so $a=\frac{7!}{8^{\frac{4}{3}}}$.
Step 7: Substitute $a$ using $r=2$ instead of current assumed, $a=\frac{7}{8}$.
Step 8: Additionally, the additional mistake in calculation keeping the factor in step 7 wrong.
Step 9: We correct this calculation by keeping initial conditions wrong and final contradiction/ error remains.
\end{proof}

\end{document}
